
\section{Youtube: All Slavic Cuisines Explained}

Video na adrese: \url{https://www.youtube.com/watch?v=zWXlwhPlyhQ}


Odpověď na otázku, proč si myslím, že trdelník původně pochází z Německa nebo Rakouska:

Z Německa pocházejí v literatuře nejstarší zmínky a recepty pro variantu, kterou dneska nazýváme trdelník nebo kurtoskalács. V primitivní verzi vyšel recept už v půlce 15. století, ale myslím, že skutečný základ pro recept, který se v následujících stoletích rozvíjel vyšel v 16. století pod názvem Einen Eyerkuchen oder Spißkuchen zu machen ve dvou kuchařkách. V anonymní Klosterkochbuch z půlky 16. století a v roce 1593 v kuchařce Buch III: Vom Kochen. Aus: Oeconomia. Oder Haußbuchs od Johannese Colera. Jednalo se o kynuté těsto s rozinkami, obarvené šafránem, omotané v pruhu kolem špízu. Při pečení se potíralo máslem a nepoužívaly se žádné posypky ani cukr. V následujících staletích vyšly různé varianty tohoto receptu v dalších německých a rakouských kuchařkách, třeba v Freywillig-auffgesprungener Granat-Apffel Ein gantz neues und nutzbahres Koch-Buch od Eleanory Lichtenstein z přelomu 17. a 18. století a na konci 18. století v maďarských kuchařkách  Márie Mikes a Némely étkek készítési módja od Kristófa Simaie (ta vyšla v Kremnici na Slovensku).

Nejstarší recept s posypkou, který se skládala z cukru a mandlí jsem našel v kuchařce od Judith Marie Rohrer, která vyšla v Brně česky i německy v roce 1831. V češtině je to recept č. 386 Stonkové neb uvinuté koblížky v kuchařce Co hrdlo ráčj, neb, 500 rozličných pokrmů, pochoutek, rosolů. Ten název je očividně překlad z Prügelkrapfen, který se používal v Rakousku (další názvy používané v Rakousku byly Spiesskrapfen a Ringelkrapfen). Nevylučuji, že kopírovala z nějaké starší rakouské kuchařky, kterou jsem zatím nenašel v digitální verzi. Hodně podobné recepty vyšly v Maďarsku v 70. letech 19. století, jeden v kuchařce od Eszter Vörös, druhý od Terez Doleskó. Ten druhý se označuje za první recept na moderní kurtoskalács, protože se těsto obalovalo v cukru už před pečením, takže se u něj vytvoří karamelová krusta. Podle mě vycházely z toho moravského receptu od Rohrer. Jinak to, že se do Maďarska špízové koláče rozšířily z Rakouska nebo Německa tvrdí maďarští autoři, kteří se zabývali historií kurtoskalácse (Ferenc Pozsony a Péter Hantz). Nejstarší zmínky z maďarských pramenů, které našli, jsou až ze 17. století. To tvrzení, že kurtoskalács vznikl ve středověku při útocích Tatarů je legenda, kterou v 19. století publikoval spisovatel Balázs Orbán. 

V Německu vzniky v 16. století i další dvě varianty špízových koláčů. První byla tvořená pruhem těsta, který byl stejně široký jako délka pečící formy, takže se nenamotával jako provaz, omotal se kolem ní v jednom kusu. Tato varianta vyšla poprvé jako Spiesskuchen v Ain künstlichs und nutzlichs Kochbuch od Balthasara Staindla v roce 1547 a později v Ein new Kochbuch od Marxe Rumpolta z roku 1581. V roce 1694 podobný recept pod názvem Bohmische Küchlein v kuchařce Haus- Feld- Arzney- Koch- Kunst und Wunderbuch od Johanna Christopha Thiema. Nepodařilo se mi zjistit, proč začal být spojovaný s českým prostředím. Böhmische Kuchlein zmiňuje i Čeněk Zíbrt ve svém Staročeském umění kuchařkém, cituje Das kleine Nürnberger Kochbuch z roku 1727, ale tu se mi nepodařilo dohledat v digitální verzi. Vrcholu tato varianta došla v Neues Saltzburgisches Koch-Buch od Conrada Haggera z první třetiny 18. století, kde je zdobená různými figurkami z marcipánu. V Kleines Kochbuch, oder Sammlung einiger wenigen, aber ganz besondern ausgesuchten Speisen, nach heutigem wienerischen Geschmacke or Marie Anny Wieser z roku 1790 je recept na Spiesskrapfen, plát těsta se peče s cukrovou polevou a sype pistáciemi, což je asi nejstarší recept na kteroukoliv variantu špízového koláče s posypkou se semínek nebo ořechů.

Poslední varianta z 16. století je z litého těsta. Dneska je známá z Německa jako Baumkuchen, v Rakousku Prügelkrapfen, v Litvě Šakotis, ve Švédsku Spettekaka, rozšířená je ale i jinde. První recept vydala pod názvem Wiltu ein spiß kuchenn bachenn Susanna Harsdörfer v kuchařce Was zu 180 Zuckern Leckkuchleinn von Wirtz und Zucker genomen wirt v roce 1582. Tato varianta byla nejpopulárnější a i dneska jsou její varianty nejrozšířenější (když nepočítáme turistické stánky s trdelníkem a kurtoskalácsem, které jsou už po celém světě).  Recept obsahuje řada německých a rakouských kuchařek z 18. století. Na přípravu je mnohem složitější než předchozí dvě a málokde se z ní stalo lidové jídlo. Většinou se připravovala ve šlechtických kuchyních a později v cukrárnách. Zajímavé je, že nejstarší český recept s názvem trdelník byl na tuto variantu. Vyšel v českém překladu kuchařky Moje skrz čtyřicetileté vykonávání známá Kuchařská kniha pro velké i menší tabule od Sibylly Dorizio, nejspíš kolem roku 1816. To ale nebylo poprvé, co v českém prostedí vyšel. Už v roce 1763 obsahovala V nově rozmnožená knížka kuchařská recept Velmi dobrá postní lahůdka na rožně, recept  na tuto litou variantu. Opět je to ale překlad z němčiny. Další recept vyšel v nejdůležitější kuchařce vyšlé v Čechách na začátku 19. století, Die Bayersche Koechin in Böhmen od Anny Marie Neudecker z roku 1805, pod názvem 558 Schliff- oder Spiesstorte. Neudecker byla autorka, která do naší kuchyně přinesla guláš nebo řízky. Pod názvem trdelník vyšel i v roce 1912 ve vydání Velké kuchařky M. D. Rettigové sestavené Eliškou Neubaurovou a Idou Caligris. Jinak ale byla na přelomu 19. a 20. století rozšířená pod názvem trdelovec. Prodávala se v Praze v cukrárnách a na různých poutích, často pod názvem Staročeský trdlovec, což bylo podobně jako dneska u Staročeského trdelníku terčem vtipů. Na Znojemsku se z této varianty stalo i lidové jídlo, hlavně v německy mluvících oblastech. Po 2. světové válce a odsunu Němců tu ale zanikla. Dodnes se ale občas vyskytuje v Rakouských obcích u českých hranic. Stejně tak po 2. světové válce vymizel trdlovec z Prahy a dneska je úplně zapomenutý.

Pokud jde o rozšíření v Českých zemích, tak už jsem zmínil několik receptů, z nichž nejdůležitější byly stonkové koblížky od Judith Marie Rohrer. To byl základ receptu tak jak ho známe dnes. Jak je vidět, používaly se různé názvy, ne jen trdelník. Nejstarší výskyt použití termínu trdelník jsem našel z roku 1724, kdy v Brně dostal advokát jako odměnu za svoje služby černou zvěřinu, ryby, při tom jeden trdelník obecní. Myslím si ale, že v tomto případě trdelník označuje trdelného, čili pohlavně dospělého, kapra. To je totiž další význam slova trdelník. Jinak se název trdelník objevuje až na začátku 19. století v několika zdrojích. Kromě receptu od Sibylly Dorizio se jedná o sbírku básní Muza Morawská od Josefa Heřmana Agapita Gallaše z roku 1813 německo-český slovník Karla Ignáce Tháma z roku 1814. Gallaš žil v Hranicích a ve svých básních popisoval život na rozmezí Hané a Valašska. Trdelník se vyskytuje ve třech básních a je pro něj jedním ze symbolů prostopášnosti a kuchyně vyšších vrstev. Takže v této době a místě to bylo spíš měšťanské nebo šlechtické jídlo. U Tháma se vyskytuje překlad dvou německých názvů používané pro špízové koláče. Rakouské označení Prügelkrapfen překládá jako trdelník, starší německé označení Spiesskuchen pak jako vaječník.Termin vaječník se vyskytuje i v dalších českých slovnících z první poloviny 19. století. například u Josefa Jungmanna jsou trdelník a vaječník synonyma a vaječník samotný popisuje jako: koláč vaječný (Ros.) jest pečité jídlo, těsto na koláče se rozválí, při plamenu na způsob kuřete se otáčí a peče, pak cukrem posypané se jedí. 

S termínem vaječník je problém v tom, že má ještě víc významů, než trdelník. Například pohlavní orgány. Ale i v kontextu jídla. Například Jungmann jako další význam tohoto slova uvádí omeletu nebo bábovku. Takže je skutečně složité ve starší literatuře odlišit, kdy se vaječníkem skutečně myslí špízový koláč. Obecně se to dá odlišit ve slovnících, kde se používá jako překlad německého slova spiesskuchen, které jiný význam, než špízový koláč nemá. To je třeba ve dvou slovnících od Jana Amose Komenského, Januae Lingvarum Reseratae Aureae Vestibulum a Janua Lingvarum Reserata Aurea, v české verzi Zlaté dveře jazykův otevřené. Ty vyšly v 17. století, v německo-latinsko-česko-maďarském vydání  prvního z nich  z roku 1722 je pak i explicitní česká poznámka, že vaječník = na trdle koláč. Jinak spiesskuchen cituje i Vokabulář latinský a český, nyní vnově spravený a rozšiřený od Vojtěcha Jiřího Koniáše z roku 1704 nebo  Vocabularium trilingue z roku 1689. Thesaurus linguae Bohemicae od  Václava Jana Rosy cituje Jungnmann, ale vaječník tu označuje jen jako vaječný koláč, což je dost obecné označení. Daniel Adam z Veleslavína: Nomenclator omnium rerum propria nomina tribus linguis – Latina, Boiemica et Germanica z roku 1586 označuje vaječník jako chléb obarvený a nasáklý šafránem a žloutkem vejce. Velmi užitečná kníška mládencóm z roku 1532 označuje vaječník jako ayerkuch, což je termín, který se použil pro špízový koláč v německé kuchařce Klosterkochbuch pocházející ze stejné doby. Ayerkuch a jeho varianty eyerkuchen nebo eierkuchen ale mají stejný problém jako vaječník, protože můžou označovat i úplně jiné pokrmy.  

Nejstarší výskyt slova vaječník jsem našel ve slovnících Mistra Klareta, které pocházejí z doby kolem roku 1360. Je tam v seznamu pečiva: kolacz, zemle, pernik, syrnik, koblih, peczen, prha, vaječník, huscze, preczlik, bochnecz. Jako latinský překlad vaječníku je použito ovarium, což znamená pohlavní orgán, což je v kontextu toho, že je tu očividně myšlený jako pečivo očividně špatně. Jestli ale označuje špízový koláč se nedá poznat, takže tento výskyt beru jako nejistý. Teoreticky to ale může znamenat, že u nás byl špízový koláč známý už v půlce 14. století, dřív, než v Německu. Ale to je jen spekulace. Další spekulativní výskyt vaječníku jako jídla ze z roku 1464, který opět cituje Čeněk Zíbrt ve Staročeském umění kuchařském. V dopise o cestě saské vévodkyně do Prahy se píše, jak jí po cestě vítali vesničané: Stará vévodkyně saská jela do Prahy. Král Jiří Poděbradský postaral se,aby všude, kudy jela od Mostu, byla vítána nadšeně, uctivě od jásajícího lidu selského. Všude ve vesnicích sedláci jí běželi naproti, nosili jí do vozu mazance, sýry, vaječníky atd. Poslední spekulativní výskyt je v Kuchařství Bavora Rodovského z Hustiřan, kde je recept na varmuži z vaječníka. Recept na vaječník samotný nepopisuje, ani nepopisuje, co to vlastně je, ale z kontextu jde o pečivo rozkrájené na kostičky a uvařené s vajíčky a jablky. V tomto případě si myslím, že skutečně jde o špízový koláč. 

Jeden potvzený výskyt špízového koláče je z Jáchymova, kde v roce 1554 kněz Johannes Mathesius v kázání kritizoval místní ženy, že rády pijí pivo a jedí spiesskucheny, místo toho, aby se staraly o své manžely. Tohle je celkem známá zmínka, kterou citují i některé české historie trdelníku, například Česká televize v pořadu What the fact, namluvená Jankem Rubešem. Samozřejmě se dá namítnout, že Jáchymov byl v tu dobu spíš německé město, takže to s českým prostředím nemá moc společného. Ale i z dalších zdrojů je jisté, že minimálně od půlky 16. století byly špízové koláče v našem prostředí známé.

Pokud jde o rozšíření trdelníku jako lidového pečiva na Moravě,  nepodařilo se mi  zjistit, kdy se vlastně rozšířily. Kromě zmínek z počátku 19. století, které už jsem citoval, se po většinu 19. století moc neobjevují. Třeba v roce 1848 se podávaly v Ivančicích na svatbě, nebo otec manželky Viléma Mrštíka jedl trdelník jako dítě kolem poloviny 19. století na předměstí Olomouce. Množství zmínek se objevuje až v 80. letech 19. století s rozvojem zájmu o lidovou kulturu a pořádáním různých vlasteneckých slavností, kde se trdelník často podával. Jeho rozšíření nejlépe dokládají národopisné výstavky, které se v té době konaly, kde byl trdelník vystavovaný na celé Moravě od Dačic přes Brno, Olomouc a pak hodně na jižní Moravě. V zásadě byl ale prezentovaný jako starodávné pečivo, které už v té době málokdo znal. Takže očividně byl rozšířený spíš v dřívější době, jen o tom nikdo nepsal, nebo je to v pramenech, které zatím nebyly digitalizované. 

Na otázce proč trdelník zmizel se čeští autoři celkem shodují, souvisí to s nahrazením černých kuchyní sporáky. To probíhalo už od 18. století, nejdřív v bohatších vrstvách, ale na konci 19. století už černé kuchyně vymizely skoro všude. To je důvod, proč existuje spousta receptů v německo jazyčných šlechtických kuchařkách z 16. až 18. století, ale v 19. století už skoro vymizel. A do půky 20. století už skoro vymizel i z lidové kuchyně. 
